\section{Conclusion}
\label{sec:conclusion}

GATE-KD treats cross-dimensional embedding distillation as an interface problem,
using a fixed PCA subspace and closed-form gauge alignment with native-space
$H_0$ regularization. Across three teacher--student configurations, this
final-embedding-only formulation obtains the numerically highest aggregate
scores among the evaluated distilled students while recording the lowest
update time and cumulative optimization time within the measured minibatch
loop. Controlled analyses show that fitted coordinate comparability provides
the dominant gain: a per-minibatch rotation-invariant control trails GATE-KD by
2.05 points, while separate schedule controls identify repeated refitting as a
secondary optimization refinement with a smaller numerical increment.
Native-space $H_0$ adds 0.53 points in the compact comparison and performs
best among the evaluated structural auxiliaries.
The current sentence-level objective does not fully preserve the teacher's
zero-shot retrieval capability, leaving broader retrieval transfer for future
work. In addition, the current corpus uses training text from three downstream
source benchmarks, so its source-task results measure label-free task adaptation
rather than corpus-disjoint transfer. The conclusions are therefore limited to
the task-adapted sentence-level regime evaluated here.
